\documentclass[11pt,a4paper]{article}
\usepackage[utf8]{inputenc}
\usepackage[margin=1in]{geometry}
\usepackage{amsmath,amssymb,amsthm}
\usepackage{listings}
\usepackage{xcolor}
\usepackage{hyperref}

\newtheorem{theorem}{Theorem}
\newtheorem{lemma}[theorem]{Lemma}

\lstset{
  language=C++,
  basicstyle=\ttfamily\small,
  keywordstyle=\color{blue}\bfseries,
  commentstyle=\color{green!60!black},
  stringstyle=\color{red!70!black},
  numbers=left,
  numberstyle=\tiny\color{gray},
  breaklines=true,
  frame=single,
  tabsize=4,
  showstringspaces=false
}

\title{IOI 2011 -- Garden}
\author{Solution}
\date{}

\begin{document}
\maketitle

\section{Problem Summary}
A garden has $N$ nodes and $M$ edges (trails), each with a beauty value (lower index = more beautiful). From each node, you always take the most beautiful available trail. Because you cannot immediately backtrack, each node effectively has a unique successor. Given a target node~$P$ and a step count~$K$ (up to $10^{18}$), count how many starting nodes reach $P$ after exactly $K$ steps.

\section{Solution}

\subsection{Functional Graph Construction}
Since there is no immediate backtracking, split each node $v$ into two states:
\begin{itemize}
    \item State $(v, 0)$: arrived via a non-best edge (or starting). Leave via the best edge.
    \item State $(v, 1)$: arrived via the best edge. Leave via the second-best edge (or the best edge if no second-best exists).
\end{itemize}
This gives a functional graph on $2N$ states, each with exactly one successor.

\subsection{Cycle Detection and Counting}
In a functional graph, every state eventually enters a unique cycle. For a target state $t$, we need to count starting states $s$ (of the form $(v, 0)$) such that $f^K(s) = t$, where $f$ is the successor function.

\textbf{Approach:}
\begin{enumerate}
    \item For each of the two target states $t \in \{(P,0), (P,1)\}$:
    \begin{enumerate}
        \item Detect the cycle containing $t$ (or the cycle $t$'s rho-path leads into).
        \item Build the reverse graph. Identify which states lie on the cycle and which are on ``tails.''
        \item For tail states: the distance to $t$ is unique and finite. Check if it equals $K$.
        \item For cycle states at distance $d$ from $t$: the state reaches $t$ after $d + k \cdot C$ steps for any $k \ge 0$, where $C$ is the cycle length. Check if $K \equiv d \pmod{C}$ and $K \ge d$. Then BFS backward through non-cycle predecessors.
    \end{enumerate}
    \item Sum the counts from both target states (no double-counting occurs since each starting state follows a unique path).
\end{enumerate}

\subsection{Complexity}
\begin{itemize}
    \item \textbf{Time:} $O(N)$ per query (amortized over the BFS).
    \item \textbf{Space:} $O(N)$.
\end{itemize}

\section{Implementation}

\begin{lstlisting}
#include <bits/stdc++.h>
using namespace std;

int main() {
    ios::sync_with_stdio(false);
    cin.tie(nullptr);

    int N, M, P, Q;
    cin >> N >> M >> P >> Q;

    vector<vector<pair<int,int>>> adj(N); // (edge_idx, neighbor)
    for (int i = 0; i < M; i++) {
        int u, v;
        cin >> u >> v;
        adj[u].push_back({i, v});
        adj[v].push_back({i, u});
    }
    for (int v = 0; v < N; v++)
        sort(adj[v].begin(), adj[v].end());

    // Build functional graph on 2N states
    // State v*2+0: go via best edge; state v*2+1: go via second-best
    int S = 2 * N;
    vector<int> nxt(S, -1);

    for (int v = 0; v < N; v++) {
        if (adj[v].empty()) continue;
        int bestEdge = adj[v][0].first;
        int bestNeigh = adj[v][0].second;
        int secNeigh = (adj[v].size() >= 2) ? adj[v][1].second : -1;
        int secEdge = (adj[v].size() >= 2) ? adj[v][1].first : -1;

        // State v*2+0: leave via best edge
        {
            int arrived = (adj[bestNeigh][0].first == bestEdge) ? 1 : 0;
            nxt[v*2+0] = bestNeigh * 2 + arrived;
        }

        // State v*2+1: leave via second-best (or best if none)
        if (secNeigh >= 0) {
            int arrived = (adj[secNeigh][0].first == secEdge) ? 1 : 0;
            nxt[v*2+1] = secNeigh * 2 + arrived;
        } else {
            int arrived = (adj[bestNeigh][0].first == bestEdge) ? 1 : 0;
            nxt[v*2+1] = bestNeigh * 2 + arrived;
        }
    }

    // Build reverse graph
    vector<vector<int>> rev(S);
    for (int s = 0; s < S; s++)
        if (nxt[s] >= 0 && nxt[s] < S)
            rev[nxt[s]].push_back(s);

    // Solve for a target state and step count K
    auto solve = [&](int target, long long K) -> int {
        if (target < 0 || target >= S || nxt[target] < 0) return 0;

        // Walk forward from target to find cycle
        vector<int> path;
        set<int> vis;
        int cur = target;
        while (!vis.count(cur)) {
            vis.insert(cur);
            path.push_back(cur);
            if (nxt[cur] < 0) break;
            cur = nxt[cur];
        }
        int cycleEntry = cur;

        // Find cycle length
        int cycleLen = 1;
        int c = nxt[cycleEntry];
        while (c != cycleEntry) { c = nxt[c]; cycleLen++; }

        // Identify cycle nodes
        set<int> onCycle;
        onCycle.insert(cycleEntry);
        c = nxt[cycleEntry];
        while (c != cycleEntry) { onCycle.insert(c); c = nxt[c]; }

        int count = 0;
        bool targetOnCycle = onCycle.count(target);

        if (targetOnCycle) {
            // Compute cycle-distance from each cycle node to target
            map<int, long long> cycleDist;
            cycleDist[target] = 0;
            int prev = -1;
            for (int r : rev[target])
                if (onCycle.count(r)) { prev = r; break; }
            long long d = 1;
            int node = prev;
            while (node != target && node >= 0) {
                cycleDist[node] = d++;
                int np = -1;
                for (int r : rev[node])
                    if (onCycle.count(r) && !cycleDist.count(r))
                        { np = r; break; }
                node = np;
            }

            // BFS backward from each cycle node through tails
            for (auto &[cn, cdist] : cycleDist) {
                if (cdist > K) continue;
                long long need = ((K - cdist) % cycleLen + cycleLen) % cycleLen;

                queue<pair<int, long long>> bfs;
                bfs.push({cn, 0});
                while (!bfs.empty()) {
                    auto [u, dd] = bfs.front(); bfs.pop();
                    long long totalDist = dd + cdist;
                    if (totalDist > K) continue;
                    if ((K - totalDist) % cycleLen == 0) {
                        if (u % 2 == 0) count++;
                    }
                    for (int r : rev[u]) {
                        if (!onCycle.count(r))
                            bfs.push({r, dd + 1});
                    }
                }
            }
        } else {
            // Target on tail: only direct predecessors at distance K
            queue<pair<int, long long>> bfs;
            bfs.push({target, 0});
            while (!bfs.empty()) {
                auto [u, dd] = bfs.front(); bfs.pop();
                if (dd == K) {
                    if (u % 2 == 0) count++;
                    continue;
                }
                if (dd > K) continue;
                for (int r : rev[u])
                    bfs.push({r, dd + 1});
            }
        }

        return count;
    };

    while (Q--) {
        long long K;
        cin >> K;
        int ans = solve(P * 2 + 0, K) + solve(P * 2 + 1, K);
        cout << ans << "\n";
    }

    return 0;
}
\end{lstlisting}

\end{document}
